% Design options (M0-3): the same slides in ten designs, light and dark.
%   \def\slidedesign{d03}  design file design-d03.sty (default d01)
%   \def\slidepalette{dark} light or dark
%   \def\slidescheme{…}     optional colour scheme: algorithmik (default), networks, graphite
% Build all: theme/designs/build_designs.sh
\ifdefined\slidedesign\else\def\slidedesign{d01}\fi
\ifdefined\slidepalette\else\def\slidepalette{dark}\fi
\ifdefined\slidescheme\else\def\slidescheme{algorithmik}\fi
\def\designbigtype{d06}
\ifx\slidedesign\designbigtype\def\designsize{14pt}\else\def\designsize{12pt}\fi
\edef\designopts{aspectratio=169,\designsize,xcolor={dvipsnames}}
\expandafter\documentclass\expandafter[\designopts]{beamer}
\usepackage{infschoolslides}
\definecolor{networksBase}{HTML}{D2601A}
\definecolor{networksDark}{HTML}{9E4310}
\IfStrEqCase{\slidescheme}{%
  {algorithmik}{\slidecolor{Rhodamine!90!black}}%
  {networks}{\slidecolor[networksDark!90!black][networksBase!80!white]{networksBase}}%
  {graphite}{\slidecolor{Rhodamine!90!black}}}
\slidesubject{Algorithmik}
\slidechapter{3}
\slidecoverimages{cover/algorithmik/playing-cards.jpg, cover/algorithmik/pencils.jpg, cover/algorithmik/maze.jpg,
  cover/algorithmik/signpost.jpg, cover/algorithmik/bridges.jpg, cover/algorithmik/subway-map.jpg}
\title{Sortieren}
\subtitle{Selectionsort · Bubblesort · Mergesort}
\author{Florian Köhl}
\makeatletter
\usepackage{designs-common}
\usepackage{design-\slidedesign}
\IfStrEq{\slidescheme}{graphite}{\dsg@graphite}{}
% titles on the page (not in a band): dark in light mode, white in dark mode
\IfStrEq{\slidepalette}{dark}{\colorlet{pageTitle}{white!95!black}}{\colorlet{pageTitle}{primaryDark}}
% sections of the deck, for designs that show the deck's progress
\newcommand{\decksections}{Einstieg, Selectionsort, Mergesort, Spannbäume}
\makeatother

\begin{document}

% ------------------------------------------------------------------ L1 Titel
{\setbeamertemplate{footline}{}\setbeamertemplate{background}{}\begin{frame}[plain]\titlepage\end{frame}}

% ------------------------------------------------------------------ the design's pitch
\section{Einstieg}
\begin{frame}[t]{\designnumber\ «\designname»}
  \setbeamerfont{itemize/enumerate body}{size=\small}\small\designidea\par\smallskip
  \begin{columns}[T,onlytextwidth]
    \begin{column}{0.47\textwidth}
      \slidelabel{Eigenschaften}
      \begin{itemize}\setlength{\itemsep}{0.5mm}\designprops\end{itemize}
      \smallskip\slidelabel{Nimm dieses Design, wenn …}
      \designpick
    \end{column}
    \begin{column}{0.47\textwidth}
      \slidelabel{Stärken}
      \begin{itemize}\setlength{\itemsep}{0.5mm}\designpros\end{itemize}
      \slidelabel{Schwächen}
      \begin{itemize}\setlength{\itemsep}{0.5mm}\designcons\end{itemize}
    \end{column}
  \end{columns}
\end{frame}

% ------------------------------------------------------------------ L2 Abschnitt
\section{Selectionsort}
\sectionslide{3.2}{Selectionsort}{«Wiederholtes Suchen des Minimums»}

% ------------------------------------------------------------------ L4 Aufzählung
\begin{frame}{Effizienz: Laufzeit}
  \begin{itemize}
    \item Wie «schnell» ist Selectionsort?
    \item Wie messen wir die Effizienz eines Algorithmus?
    \begin{itemize}
      \item Zeitmessung? Hängt von Hardware und Implementierung ab.
      \item Alternativen?
    \end{itemize}
    \item Grobe Idee: die \term{Anzahl Operationen} zählen
    \begin{itemize}
      \item Vergleiche zwischen zwei Elementen
      \item Vertauschungen zweier Elemente
    \end{itemize}
  \end{itemize}
\end{frame}

% ------------------------------------------------------------------ L7 Begriff
\begin{frame}{Sortiert?}
  \begin{concept}[Aufsteigend sortiert]
    Eine Liste ist \term{aufsteigend sortiert}, wenn jedes Element kleiner oder gleich gross ist
    wie alle nachfolgenden Elemente.
  \end{concept}
  \begin{itemize}
    \item Sind diese Karten sortiert?
    \item Und die leere Liste?
  \end{itemize}
\end{frame}

% ------------------------------------------------------------------ L6 Grafik (cards)
\begin{frame}{Selectionsort in Aktion}
  \vspace*{2mm}
  \cardArray{name=c, items={2h,6h,5h,9h,7h,10h}, parts={1/1/S,2/6/U}, size=0.12}
  \swaparrow{c2}{c3}\pointerbelow{c3}{min}

  \vspace{12mm}\centering Das kleinste Element in U ist die 5: Wir vertauschen es mit dem ersten Element von U.
\end{frame}

% ------------------------------------------------------------------ L10 Code
\begin{frame}[fragile]{Selectionsort in Python}
  \begin{code}[selection\_sort.py]
def selection_sort(liste):
    for i in range(len(liste) - 1):
        k = i  # Position des Minimums in U
        for j in range(i + 1, len(liste)):
            if liste[j] < liste[k]:
                k = j
        liste[i], liste[k] = liste[k], liste[i]
  \end{code}
\end{frame}

% ------------------------------------------------------------------ L13 Aufgabe, «Mach mit»
\begin{frame}[doit={zu zweit · 15 Min}]{Ihr seid dran!}
  \vfill
  \begin{task}[Sortieren von Hand]{38–41}
    \large Sortiere die Listen im Skript (S. 48) mit Selectionsort und Bubblesort von Hand.
  \end{task}
  \vfill
\end{frame}

% ------------------------------------------------------------------ L8 Algorithmus + figure
\section{Mergesort}
\begin{frame}{Mergesort: Merge}
  \begin{columns}[T,onlytextwidth]
    \begin{column}{0.5\textwidth}
      \begin{enumerate}
        \item Vergleiche die vordersten Elemente beider Listen.
        \item \alert{Füge das kleinere hinten an die neue Liste an.}
        \item Ist eine Liste leer, hänge den Rest der anderen an.
      \end{enumerate}
    \end{column}
    \begin{column}{0.46\textwidth}
      \pythonArray{items={2,5,3,9}, parts={0/2/links,2/4/rechts}, done={0}, highlight={1,2}, ptrs={1/i,2/j}}
      \vspace{4mm}
      \pythonArray{items={2,~,~,~}, blank, done={0}, drawindexes=false}
    \end{column}
  \end{columns}
\end{frame}

% ------------------------------------------------------------------ L9 Schritt für Schritt
\section{Spannbäume}
\begin{frame}{Kruskal in Aktion}
  \centering
  \begin{graph}[x=15mm, y=15mm]
    \node[vertex] (C) at (-1.2,0.4) {C}; \node[vertex] (E) at (0,1.2) {E};
    \node[vertex] (D) at (1.2,0.4) {D};  \node[vertex] (A) at (-0.5,-0.6) {A};
    \node[vertex_active] (B) at (0.8,-1) {B};
    \path (C) edge["3", onslide=<2->{chosen}] (E) (E) edge["10"] (D) (C) edge["2", chosen] (A)
          (A) edge["4", onslide=<3->{removed}] (E) (A) edge["6"] (D) (D) edge["5"] (B) (A) edge["5"] (B);
  \end{graph}

  \vspace{3mm}
  \only<1>{C–A (2) bildet keinen Kreis: \term{gewählt}.}%
  \only<2>{C–E (3) bildet keinen Kreis: \term{gewählt}.}%
  \only<3>{A–E (4) würde einen Kreis bilden: \term{verworfen}.}

  \vfill\stepdots{3}
\end{frame}

% ------------------------------------------------------------------ «Zum Lesen»
\begin{frame}[read={Skript S. 92}]{Ein Baum ist …}
  \begin{enumerate}
    \item ein \term{zusammenhängender} Graph ohne Kreise.
    \item ein Graph, in dem es zwischen zwei Knoten genau \term{einen Pfad} gibt.
    \item ein zusammenhängender Graph mit $n$ Knoten und $n-1$ Kanten.
  \end{enumerate}
  \medskip Alle drei Beschreibungen sind gleichwertig.
\end{frame}

\end{document}
